\section{Delooping groups}
\label{delooping}

\paragraph{The external point of view}
A \emph{group} consists of a set $A$, together with an operation
$\gm:A\to A\to A$ (the \emph{multiplication}), an element $\ge:A$ (the
\emph{unit}), and an operation $\gi:A\to A$ (the \emph{inverse}) such that
multiplication is associative, admits $\ge$ as unit, and $\gi(x)$ is the
two-sided inverse of any element $x:A$. We write $\Group$ for the type of all
groups, and $G\toGroup H$ for the type of group morphisms between groups~$G$
and~$H$.
% :
% \begin{align*}
  % \Group\defd
  % \Sigma(A:\Set).&
  % (\gm:A\to A\to A)\times(\ge:A)\times(\gi:A\to A)\times
  % \\
  % &((x,y,z:A)\to\gm(\gm(x,y),z)=\gm(x,\gm(y,z))\times
  % \\
  % &((x:A)\to\gm(\ge,x)=\gm(x,\ge)=x)\times
  % \\
  % &((x:A)\to\gm(\gi(x),x)=\gm(x,\gi(x))=e)
% \end{align*}
%
In the following, we use the traditional notations for groups: given two
elements $x,y:G$, we simply write $xy$ instead of $\gm(x,y)$, $1$ instead of
$\ge$, and $x^{-1}$ instead of $\gi(x)$.

\paragraph{The internal point of view}
\label{internal-groups} 
A \emph{pointed} type consists of a type~$A$ together with a distinguished
element, often written~$\pt{}$ and sometimes left implicit. Given a pointed type
$(A,\pt{})$, its \emph{loop space}~$\Loop A$ is defined as the type of paths
from $\pt{}$ to itself: $\Loop A\defd(\pt{}=\pt{})$. The elements of this type
are called \emph{loops}. By path induction one can construct, for every two
paths $p:a=b$ and $q:b=c$, a path in $a=c$ called their \emph{concatenation} and
written $p\pcomp q$. Similarly, every path~$p:a=b$, admits an \emph{inverse}
path $\sym p:b=a$. When~$A$ is a pointed groupoid, $\Loop A$ is a set, and these
operations canonically equip this set with a structure of
group~\cite[Section~2.1]{hottbook}.

\paragraph{Delooping groups}
A \emph{delooping} of a group~$G$ is a pointed connected groupoid $\B G$
together with an identification $\Beq G:\Loop\B G=G$ (we recall that a type~$A$
is \emph{connected} when the type~$\strunc{A}$ is contractible, \ie $A$ has one
connected component).
%
The notation is justified by the fact that deloopings are unique.
%
For instance, it is known that the circle is a delooping of~$\Z$: indeed, $\S1$ is a connected groupoid, and its fundamental group is~$\Z$~\cite[Section~8.1]{hottbook}.

% Local Variables:
% TeX-master: "article"
% End:
